%; whizzy paragraph -pdf xpdf -latex ./whizzypdfptex.sh
%; whizzy-paragraph "^\\\\begin{frame}"
% latex beamer presentation.
% platex, latex-beamer でコンパイルすることを想定。 

%     Tokyo Debian Meeting resources
%     Copyright (C) 2009 Junichi Uekawa

%     This program is free software; you can redistribute it and/or modify
%     it under the terms of the GNU General Public License as published by
%     the Free Software Foundation; either version 2 of the License, or
%     (at your option) any later version.

%     This program is distributed in the hope that it will be useful,
%     but WITHOUT ANY WARRANTY; without even the implied warranty of
%     MERCHANTABILITY or FITNESS FOR A PARTICULAR PURPOSE.  See the
%     GNU General Public License for more details.

%     You should have received a copy of the GNU General Public License
%     along with this program; if not, write to the Free Software
%     Foundation, Inc., 51 Franklin St, Fifth Floor, Boston, MA  02110-1301 USA

\documentclass[cjk,dvipdfmx,12pt]{beamer}
\usetheme{Tokyo}
\usepackage{monthlypresentation}

%  preview (shell-command (concat "evince " (replace-regexp-in-string "tex$" "pdf"(buffer-file-name)) "&"))
%  presentation (shell-command (concat "xpdf -fullscreen " (replace-regexp-in-string "tex$" "pdf"(buffer-file-name)) "&"))
%  presentation (shell-command (concat "evince " (replace-regexp-in-string "tex$" "pdf"(buffer-file-name)) "&"))

%http://www.naney.org/diki/dk/hyperref.html
%日本語EUC系環境の時
\AtBeginDvi{\special{pdf:tounicode EUC-UCS2}}
%シフトJIS系環境の時
%\AtBeginDvi{\special{pdf:tounicode 90ms-RKSJ-UCS2}}

\title{Erlang コードを \\ Debian パッケージにしてみた}
\author{まえだ こうへい mkouhei@debian.or.jp\\IRC nick: mkouhei}
\date{2009年5月16日}
\logo{\includegraphics[width=8cm]{image200607/openlogo-light.eps}}

\begin{document}

\frame{\titlepage{}}

\section{入門のきっかけ}
\begin{frame}[containsverbatim]{入門のきっかけ}
友人にCouchDBというデータベースを教えてもらったのがきっかけ。

\begin{enumerate}
 \item ドキュメント指向
 \item スキーマレス
 \item RESTful
 \item JSON形式でデータをやりとり
 \item ErlangとJavaScriptで実装されている
\end{enumerate}
機能的に面白いので、友人と目下活動中。少しずつ仲間を増やしています。

\end{frame}

\section{CouchDBの中身を知りたい}
\begin{frame}[containsverbatim]{CouchDBの中身を知りたい}
 \begin{center}
  \begin{enumerate}
   \item CouchDBに興味持った。
   \item CouchDBの中身を知りたい。
   \item でもErlang読めない。
   \item Erlangって並行処理とか分散処理に向いているらしい。
   \item Erlang自体も面白そうだ。
   \item じゃ、両方やるか。まずはErlangから。
  \end{enumerate}
 \end{center}
ということで、始めてみた。
\end{frame}

\section{Erlangの歴史}
\begin{frame}{Erlangの歴史}

 \begin{enumerate}
  \item エリクソンによって設計された。
  \item 分散環境、耐障害、ある程度のリアルタイム性、無停止可動での条件の
	システムを構築するためのプログラミング言語。
  \item エリクソン社内利用 → 1998年 オープンソースソフトウェアとして公
	開。
  \item Erlang Public License という Mozilla Public Licenseの派生ライセ
	ンス。
  \item 通信機器分野での実績、20年を超える歴史。
 \end{enumerate}
\end{frame}

\section{仕様}
\begin{frame}{仕様}

気になった主な仕様。

 \begin{itemize}
  \item 評価できるものは全て式
  \item 変数の書き換えはできない
  \item 関数型プログラミング言語
  \item 並行処理（並行プログラミング）
 \end{itemize}
\end{frame}

\begin{frame}{Erlangによる実装例}

代表的なもの。Debianパッケージにもなっている。

 \begin{enumerate}
  \item CouchDB
  \item ejabberd\\
	Twitterでも使われているらしい、IMサーバ
  \item YAWS(Yet another web server)
 \end{enumerate}
\end{frame}


\begin{frame}[containsverbatim]{環境を整える}

  Debianではすでにパッケージになっている。

  \begin{commandline}
   $ sudo apt-get install erlang-base erlang-mode
  \end{commandline}

\end{frame}

\begin{frame}[containsverbatim]{対話型で使ってみる}
erlと言うコマンドで、Erlang Shell(Eshell)を起動します。

 \begin{commandline}
  $ erl
 \end{commandline}

 \begin{itemize}
  \item Eshellで実行できるのは、式だけです。
  \item 式は、ピリオド``.''で終了します。
  \item 変数は英大文字から始まります。
  \item Eshellを終了するには、Ctl + C, aか、q().を実行します。
 \end{itemize}

\end{frame}

\begin{frame}[containsverbatim]{単一代入変数}
例の書き換えられない変数、というやつです。
 \begin{commandline}
  1> A=10.
  10
  2> A=10.
  10
  3> A=11.
  ** exception error: no match of right hand side value 11
 \end{commandline}
\end{frame}

\begin{frame}[containsverbatim]{イコール''=''の役割}
イコール''=''は、パターン照合を意味します。\\
未代入の変数の時だけ、値の代入を行えます。
 \begin{commandline}
  4> A1=20.
  20
  5> A2=30.
  30
  6> A3=A1+A2.
  50
  7> A3=50.
  50
  8> A3=100-50.
  50
  9> A3=A1-A2.
  ** exception error: no match of right hand side value -10
 \end{commandline}
\end{frame}

\begin{frame}{プログラムを実行する}
プログラムを実行する方法は3種類あります。
 \begin{itemize}
  \item Eshellからソースコードをコンパイルする。
  \item シェルからコマンドでコンパイルする。
  \item escriptとして実行する。
 \end{itemize}
 コンパイル自体は、ソースコードをコンパイルして生成されるBEAMファイルを
 Erlang VMが実行します。
\end{frame}

\begin{frame}{Hello Worldで試してみよう}

 
\end{frame}

\begin{frame}{現状のErlangパッケージ}
 \begin{itemize}
  \item 例えば、erlang-baseのソースパッケージには、実はBEAMしか含まれていないコードものもあり
	ます。
  \item BEAMファイルとErlang VMがあればどのアーキテクチャ上でも実行でき
	るので、先月のJava ポリシーを読んでみるのネタとも共通点がありま
	す。
  \item Erlangにはコード自動ロード機能というものがあります。
  \item 通常は、Erlang関係のパッケージは、
	/usr/lib/erlang/lib/?????/ebin/ディレクトリ以下にBEAMファイルが
	配置されることが多いです。
  \item ただし、そこにおけば、Erlangのコード自動ロード機能が働くわけではありません。
 \end{itemize}
\end{frame}

\begin{frame}{自分のコードをdebパッケージにするには}

 \begin{itemize}
  \item Makefileを用意します。
  \item BEAMファイルは、実行権限は必要ありません。
  \item debian/controlのBuild-Dependsセクションにerlang-devを、Dependsセ
	クションに\$\{erlang-base:Depends\}を追記する。
  \item CDBSで行った場合、rulesにDEB\_FIXPERMS\_EXCLUDEで、BEAMファイル
	を指定するの忘れないようにする。
 \end{itemize}

\end{frame}

\begin{frame}{debパッケージを作ってみる}


\end{frame}

\begin{frame}{まとめと今月のヨメHack}
 \begin{itemize}
  \item まだまだ、Erlang自体の勉強が必要。
  \item Perlのようにdh-erlangパッケージとか必要？
  \item ヨメに新型インフルで、スペインに行くなとかいわれているのが、面倒
	くさい。会社も自粛とか海外渡航報告とかしる、とか言ってるし。みん
	な騒ぎすぎ。
 \end{itemize}

\end{frame}

\end{document}

;;; Local Variables: ***
;;; outline-regexp: "\\([ 	]*\\\\\\(documentstyle\\|documentclass\\|emtext\\|section\\|begin{frame}\\)\\*?[ 	]*[[{]\\|[]+\\)" ***
;;; End: ***
